% appendix: prompts
\section{Prompts and decoding}
\label{app:prompts}

Every prompt was designed on validation data only and frozen in
\texttt{PREREGISTRATION.md} before any test split was opened.

\subsection{Structured-output prompt}
\TODO{verbatim prompt from Phase 5.1, together with the JSON schema shown to the model}

\subsection{Plain-answer prompt (structured output disabled)}

The structured-output cost in Section~\ref{sec:ablations} compares the same weights under two
elicitations. The plain arm uses the prompt from the Qwen3-VL technical report's own evaluation
appendix, so that this arm is anchored to the elicitation that produced the published
\result{79.1}:

\begin{quote}\ttfamily
\{question\}\\
Answer the question using a single word or phrase.
\end{quote}

\subsection{Decoding}
\TODO{temperature, top-p, max new tokens, retry policy --- from the frozen pre-registration}

\subsection{Answer normalisation}

The canonical relaxed-accuracy metric compares a gold answer of \texttt{"0"} as a \emph{string}
(Section~\ref{sec:setup:metrics}), so a prediction of \texttt{"0.0"} scores incorrect. The
normaliser therefore emits bare \texttt{"0"}. Since a measurable fraction of human-sourced ChartQA
gold tables contain a zero, this is a routine case rather than an edge case.

\TODO{full normaliser rules, frozen}
